\chapter{歴史的注}

非ユークリッド幾何学の歴史は、すでに非常によく、しばしば書かれているので、ここでは簡単な概略だけを述べる。

ユークリッドの公理のなかに、表現がいくぶん複雑で、他のもののような単純で初歩的な事実とは思えないものが一つある。それは次のものである。\footnote{ポール・タンヌリーによる、ユークリッドの公理についての論文を見よ。\emph{Bulletin des Sciences Mathématiques}, 1884.}

\begin{quotation}\emph{二直線が第三の直線で切られ、その切る線の同じ側の内角の和が二直角より小さいならば、二直線は十分に延長すれば、その側で出会う。}\end{quotation}

多くの数学者、とりわけルジャンドルが、この命題の証明を与えようと試みた。すなわち、それがそれに先立つより単純な公理の必然的な帰結であることを示そうとした。ルジャンドルは、三角形の角の和は決して二直角を超えないこと、また、この和が二直角に等しい三角形が一つでもあれば、すべての三角形についてそうであることを証明した。これはもちろん、直線が無限の長さをもつという仮定のもとでのことであった。しかし彼は、角の和が二直角であるような三角形が存在することを証明できなかった。\footnote{たとえば、彼の\emph{幾何学原論}（\emph{Éléments de Géométrie}）の第12版、第I巻、命題XIX、および注II を見よ。また、\emph{Bulletin des Sciences Mathématiques} 第一シリーズ第二巻の、非ユークリッド幾何学についての論文におけるクラインの言明も見よ。}

ついに、この言明は証明できないものであり、それがつねに真であるとは限らないと仮定しても、同じように無矛盾な幾何学を築き上げることができる、さらにはユークリッドのすべての公準は、われわれの経験によって真と受け入れるに至った仮説にすぎず、論理的な幾何学の展開においては反対の言明で置き換えうるものである、と考える数学者が現れはじめた。

この理論の端緒は、ときにガウスに帰せられてきたが、いまでは、1766年にランベルト\footnote{\emph{American Mathematical Monthly}, 1895年7--8月号を見よ。}が、平行線公理には証明が必要であると主張し、この公理が成り立たない幾何学の特徴のいくつかを与えた論文を書いていたことが知られている。さらに1733年という昔にも、イタリア人サッケーリが一冊の本を出版し、そこで非ユークリッド幾何学の完全な体系を与え、そのうえで、これらの他の仮説は偽であると独断的に断言することによって、自分自身とその本を救った。本書の第一章の基礎としてとられているのは、彼の取り扱いの方法である。\footnote{ハルステッドによるサッケーリの翻訳は、\emph{American Mathematical Monthly} に連載されてきた。}

ガウスは長年、平行線の公理を証明しようと努めていた。この公理を否定した帰結である定理のいくつかを発見していたかもしれないが、この主題について何も公表しなかった。

ロシアのロバチェフスキーとハンガリーのヨハン・ボヤイが、はじめて、平行線の公理が必ずしも真ではないことを主張し、証明した。二人の仕事は互いにまったく独立であり、それぞれがこの発見の全面的な功績を有する。彼らの結果は1830年ごろに公表された。

これらの発見が多くの注目を集めるまでには、長い時間がかかった。そのあいだに、後にわれわれの主題に多くの光を投げかけることになる、他の方面の研究が進められていた。それは説明としてではなく、顕著な類似によってであった。

たとえば、一、二年の間隔をおいて、同じ雑誌（クレレ）に、ロバチェフスキーの研究結果を述べた論文と、ミンディングの曲面についての論文が現れた。ミンディングは、$a$ などの代わりに $ia$ とおけば、球面三角法の公式が成り立つような曲面を見いだしたのである。しかし、これら二つの論文は、それらの関連が（ベルトラミによって）気づかれるまでに、三十年間も公表されたままであった。

またケイリーは、1859年に\emph{Philosophical Transactions} に「第六論文（量子論）」を発表し、そこで測定の射影的理論を展開して、あらゆる図形が、彼が絶対形と呼んだある特別な図形との間でもつ非調和関係を考えることにより、計量的性質を射影的に扱いうることを示した。1872年にクラインがこの理論を取り上げ、それが非ユークリッド幾何学の完全な像を与えることを示した。

また、平面や空間のなかを動き回るにつれて、ある法則に従って変化する測定の単位を考えれば、非ユークリッド幾何学が得られることも示されている。

これらの分野の古い研究者たちは、鋭角の仮説をとる幾何学だけを発見した。直線が有限の長さをもつという仮定を調べることは、彼らには思い浮かばなかった。楕円幾何学は、1854年に幾何学の基礎の研究に取り組んだリーマンによって発見されるまで、残されていた。彼はそれを、まったく異なる観点、抽象的な代数的観点から研究した。われわれの空間や幾何学的図形は、例示のためにあげるにとどめて、変数の系を考えたのである。彼は、これらの変数の関数で、長さあるいは距離の要素と呼べるものの性質はどのようなものか、という問題を調べ、最も簡単な場合には、それは、変数の微分の二次式の平方根でなければならず、その係数はそれ自身変数の関数でありうることを見いだした。二次式の異なる形をとることにより、これらの異なる種類の幾何学が無数に得られるが、そのほとんどでは、物体は大きさも形も変えずに動かせるという公理を失う。

この概観には、さらに二つの名前を加えるべきである。ヘルムホルツとクリフォードである。彼らは科学雑誌に論文を発表して、この主題を広く知らしめるのに大いに貢献した。楕円空間における平行線の理論は、前に説明したとおり、クリフォードに負っている。彼は、この幾何学のなかに、ユークリッド幾何学が成り立つ有限な曲面がありうることを示した。\footnote{非ユークリッド幾何学のより興味深い解説のいくつかは次のとおりである。\emph{Encyclopedia Britannica}、サー・ロバート・ボールによる項目「測定」（``Measurement''）。\emph{Revue Générale des Sciences}, 1891、ポアンカレによる「非ユークリッド幾何学」（``Les Géométries Non-Euclidean''）。\emph{Bulletin of the American Mathematical Society}, 1900年5月および6月、フレデリック・S・ウッズによる「ロバチェフスキーの幾何学」（``Lobachevsky's \emph{Geometry}''）。\emph{Mathematische Annalen}, Bd.~xlix, p.~149, 1897、および\emph{Bulletin des Sciences Mathématiques}, 1897、「ガウスとボヤイの書簡」（``Letters of Gauss and Bolyai''）。とくに興味深いのは、互いに平行な三本の直線が、つねに同じ有限の面積を囲むことができるという仮説のもとでの三角形の面積の公式を、ガウスが与えている一通の手紙である。最後の二つの論文は、エンゲル教授とシュテッケル教授の出版物に言及しており、それらはドイツ語で、平行線の理論の完全な歴史と、ロバチェフスキーおよびボヤイの著作と生涯を与えている。また、ジョージ・ブルース・ハルステッド教授による、ロバチェフスキーとボヤイの翻訳、およびヴァシリエフ教授の講演の翻訳も見よ。}

非ユークリッド幾何学の主な教訓は、幾何学の公理は、物理科学の理論のように、われわれの経験からの演繹にすぎないということである。数学者にとって、それらは真偽が彼の関知するところではなく、哲学者のみに関わる仮説である。彼はそれらを望む形でとり、その上に幾何学を築いてよく、こうして得られた幾何学は、数学の他の分野に応用をもつ。

これらの幾何学的命題に適用されるかぎり、「公理」は「自明」ではなく、必ずしも真ではない。ある言明が証明できる、すなわち、すでに採用された公理の必然的帰結であるならば、それを公理と呼ぶべきではない。互いに矛盾する二つ以上の言明が、すでに受け入れた公理のすべてと同じように無矛盾であるときには、そのどちらをとってもよく、われわれが選んだ言明が、われわれの幾何学にとっての公理となる。われわれの幾何学とは、この公理の帰結の研究である。

これまで研究してきた三種類の幾何学を区別する仮定は、異なる形で表現できる。点を通って、一本、二本、あるいは零本の平行線が引けると言うこともできる。あるいは、三角形の角の和が二直角に等しいか、より小さいか、より大きいと言うこともできる。あるいは、直線が無限遠に二つの実点、一つの実点、あるいは実点をもたないと言うこともできる。あるいは、平面上で相似な図形がありうるか、ありえないか、直線の長さが有限か無限か、などと言うこともできる。しかしこれらのどの形も幾何学の性質を決定し、他のものはそこから導かれる。
